\subsection{GraphQL API}
\label{subsec:github_api_graphql}
GitHub poskytuje také GraphQL \API, které klientům umožňuje dotazovat se přesně na požadovaná data. Jedná se o~odlišný přístup k~získávání dat, který přináší vyšší flexibilitu a~často také menší objem přenášených dat. Zatímco \REST \API vrací předem definované objekty z~více endpointů, GraphQL využívá jeden endpoint, přes který je možné vyžádat přesně ta data, která jsou potřeba~\cite{github_graphql_api,github_rest_api}.

Datový dotazovací jazyk GraphQL má následující vlastnosti:
\begin{description}
  \item[Specifikace] Na \API serveru je definováno schéma, podle kterého jsou validovány klientské dotazy.
  \item[Silná typovost] Ve schématu jsou přesně definovány typy a~vazby mezi objekty.
  \item[Introspektivita] Je možné získat podrobné informace o~schématu.
  \item[Hierarchická struktura] Tvar GraphQL dotazu odpovídá tvaru \JSON dat, která jsou vrácena.
  \item[Aplikační vrstva] GraphQL není pouze modelem uložení dat ani databázovým dotazovacím jazykem. Označení \textit{Graph} odkazuje na grafové struktury definované ve schématu, kde uzly představují objekty a~hrany vztahy mezi nimi.
\end{description}

Mezi časté problémy \REST \API patří nadměrné stahování dat (over-fetching) a~nedostatečné stahování dat (under-fetching)~\cite{github_graphql_api,github_rest_api}.

\begin{description}
  \item[Over-fetching] nastává, když klientská aplikace potřebuje pouze část vlastností objektu, ale v~\REST \API je vrácen celý objekt s~mnoha dalšími daty. To vede k~neefektivnímu přenosu dat.
  \item[Under-fetching] nastává, když klientská aplikace potřebuje více dat, než je dostupné v~jedné odpovědi \REST \API, a~z~toho důvodu musí být provedeno více požadavků na různé endpointy.
\end{description}

GraphQL \API tyto problémy omezuje tím, že umožňuje získat pouze ta data, která jsou skutečně potřebná. Tím může být dosaženo výrazného snížení přeneseného objemu dat. Současně je však třeba počítat s~tím, že dotazy mohou být velmi komplexní a~náročné na zpracování~\cite{github_graphql_api}.

Dotazy jsou v~GraphQL obvykle odesílány pomocí \HTTP metody \code{POST} na jediný endpoint. Operace, které mají být provedeny, i~požadovaná návratová data jsou definovány v~těle požadavku~\cite{github_graphql_api}. Ukázka dotazu a~odpovědi je uvedena ve \customref{Výpisech}{lst:graphql_request_example} a~\ref{lst:graphql_response_example}.

\begin{listing}[H]
\begin{minted}{graphql}
query {
  repository(owner: "rust-lang", name: "rust") {
    issues(first: 3, states: OPEN) {
      nodes {
        number
        title
        createdAt
        labels(first: 5) {
          nodes { name }
        }
      }
    }
  }
}
\end{minted}
  \caption{Ukázka GraphQL dotazu na otevřené issues repozitáře.}
  \label{lst:graphql_request_example}
\end{listing}

\begin{listing}[H]
\begin{minted}{json}
{
  "data": {
    "repository": {
      "issues": {
        "nodes": [
          {
            "number": 138540,
            "title": "Fix ICE in borrow checker",
            "createdAt": "2026-04-20T10:32:00Z",
            "labels": {
              "nodes": [
                { "name": "C-bug" },
                { "name": "T-compiler" }
              ]
            }
          }
        ]
      }
    }
  }
}
\end{minted}
  \caption{Zkrácená odpověď na GraphQL dotaz z~\customref{Výpisu}{lst:graphql_request_example}.}
  \label{lst:graphql_response_example}
\end{listing}

Rozlišovány jsou dva hlavní typy operací: \textit{Query} pro získávání dat a~\textit{Mutation} pro změnu dat. Operace typu \textit{Query} jsou určeny výhradně pro čtení dat a~umožňují získat komplexní strom informací v~jednom dotazu. Operace typu \textit{Mutation} jsou určeny pro vytváření, úpravu i~mazání dat. Používány jsou například při vytvoření nové issue, změně stavu \PR nebo odstranění komentáře~\cite{github_graphql_api}.

Autentizace pro GraphQL \API funguje obdobně jako pro \REST \API, tedy pomocí autentizačního tokenu v~hlavičce požadavku~\cite{github_graphql_api,github_rest_api}. Na GitHubu je pro GraphQL endpoint vyžadována autentizace, protože musí být řízena náročnost dotazů a~spotřeba prostředků~\cite{github_graphql_api}.
